\section{Experimental Design}
\label{sec:experimental-design}
%
% PURPOSE OF THIS CHAPTER
% Freeze the hypotheses, generated benchmark matrix, compared models, tuning
% rules, and statistical protocol before presenting results.
%
% \subsection{Research hypotheses}
% - Formulate testable hypotheses about horizon, lag distance, nonlinear
%   operators, interactions, irrelevant variables, noise, and mechanism
%   complexity.
% - Separate forecasting hypotheses from mechanism-recovery and explanation
%   hypotheses.
%
% \subsection{Synthetic mechanism families and difficulty axes}
% - Define controlled easy/medium/hard strata using measurable generator
%   properties rather than forecast error observed after the fact.
% - Cover autoregression, cross-variable lags, interactions, nonlinearities,
%   temporal scales, irrelevant variables, stochastic scale, and any implemented
%   regime changes.
% - Include random-coverage generation in addition to controlled factorial
%   subsets and report the resulting distribution of formulas.
%
% \subsection{Dataset construction}
% - Report formula counts, trajectories per formula, lengths, context, horizons,
%   burn-in, train/validation/test construction, and seed policy.
% - Prevent formula or trajectory leakage across evaluation partitions.
% - State rejection rates and reasons for discarded candidate mechanisms.
%
% \subsection{Compared models and information received}
% - Chronos-2 and Toto 2.0 as the prioritized pretrained forecasters.
% - Naive and seasonal-naive baselines needed to contextualize accuracy.
% - AutoARIMA/ARIMAX under explicitly separated target-only and exogenous-input
%   settings.
% - An explainable-by-design or symbolic system-identification model when the
%   final implementation and time budget permit.
% - Add a table showing tensor/dataframe shape, variables, lags, covariates,
%   horizon handling, fitting, and explanation access for every model.
%
% \subsection{Explanation methods and comparison protocol}
% - Match gradient with generator gradient and SHAP-like output with the declared
%   generator Shapley game.
% - Define baselines, backgrounds, coalition samples, random repetitions, and
%   estimator budgets.
% - State how intrinsic model outputs are converted to the canonical source-lag-
%   horizon representation.
%
% \subsection{Statistical analysis and ablations}
% - Specify repeated seeds, confidence intervals, paired comparisons, multiple-
%   comparison control if used, and aggregation levels.
% - Plan ablations for operator sets, lag bands, noise, number of mechanisms,
%   trajectories per mechanism, and explanation sampling budget.
% - Report both overall performance and stratified failure analysis.
%
% \subsection{Hardware and runtime protocol}
% - Record hardware, software versions, precision, batch size, worker count,
%   caching, warm-up, and timing method.
% - Compare serial CPU, parallel CPU, and GPU paths only where the implementation
%   offers equivalent semantics.
